% Source: https://artofproblemsolving.com/community/c3h2128742_trigonometry_question
% Stable identity: aops-topic-2128742-q01
Use the usual notation, with $a$, $b$, $c$ opposite $A$, $B$, $C$.

\begin{asy}
usepackage("amsmath");
size(4.5cm);
defaultpen(fontsize(10pt));
pair A=(0,0), B=(5.2,0), C=(1.7,3.4);
draw(A--B--C--cycle,linewidth(0.9));
dot(A); dot(B); dot(C);
label("$A$",A,SW); label("$B$",B,SE); label("$C$",C,N);
label("$c$",(A+B)/2,S); label("$a$",(B+C)/2,NE); label("$b$",(A+C)/2,NW);
draw(arc(A,0.55,0,degrees(C-A)),gray); draw(arc(B,0.55,degrees(C-B),180),gray); draw(arc(C,0.5,degrees(A-C),degrees(B-C)),gray);
label("\mbox{}",(min(currentpicture,user=true)+max(currentpicture,user=true))/2,black); // keeps the text after the figure black when LaTeX typesets the labels inline
\end{asy}

By the law of cosines, $\cos A=\dfrac{b^2+c^2-a^2}{2bc}$, $\cos B=\dfrac{c^2+a^2-b^2}{2ca}$ and $\cos C=\dfrac{a^2+b^2-c^2}{2ab}$. Hence
\[ab^2\cos A=\dfrac{abc}{2}\cdot\dfrac{b^2+c^2-a^2}{c^2},\qquad bc^2\cos B=\dfrac{abc}{2}\cdot\dfrac{c^2+a^2-b^2}{a^2},\qquad ca^2\cos C=\dfrac{abc}{2}\cdot\dfrac{a^2+b^2-c^2}{b^2}.\]
Dividing the given equalities by $\dfrac{abc}{2}>0$ and writing $X=a^2$, $Y=b^2$, $Z=c^2$, there is a number $k$ with
\[Y+Z-X=kZ,\qquad Z+X-Y=kX,\qquad X+Y-Z=kY .\]
Adding these three equations gives $X+Y+Z=k(X+Y+Z)$, and $X+Y+Z>0$, so $k=1$. Then the first equation gives $Y=X$ and the second gives $Z=Y$. So $a^2=b^2=c^2$, hence $a=b=c$ and the triangle is equilateral. $\blacksquare$
